\documentclass[a4paper,11pt]{article}

\usepackage[utf8]{inputenc}
\usepackage[T1]{fontenc}
\usepackage[ngerman]{babel}
\usepackage{amsmath}
\usepackage{amssymb}
\usepackage{xcolor}
\usepackage{colortbl}
\usepackage{tikz}
\usepackage{geometry}
\geometry{a4paper,margin=2.2cm}

% ----- Dark Mode (Pflicht) -----
\definecolor{bgdark}{HTML}{1E1E1E}   % dunkler Hintergrund
\definecolor{fgtext}{HTML}{E6E6E6}   % heller Text
\definecolor{accent}{HTML}{89B4FA}   % Akzent (hellblau)
\definecolor{accent2}{HTML}{F9E2AF}  % Akzent (gelb)
\definecolor{gridcol}{HTML}{4A4A4A}  % gedämpfte Gitterfarbe
\definecolor{rulecol}{HTML}{6A6A6A}  % helle Linienfarbe
\pagecolor{bgdark}
\color{fgtext}

% helle Linien
\setlength{\arrayrulewidth}{0.3pt}
\arrayrulecolor{rulecol}

% ----- Karo-Raster zum Ausfüllen -----
% #1 = Höhe des Karobereichs (z. B. 5cm)
\newcommand{\karo}[1]{\par\noindent
  \begin{tikzpicture}
    \draw[step=5mm, gridcol, very thin] (0,0) grid (\linewidth,#1);
  \end{tikzpicture}\par}

% ----- Kopf -----
\newcommand{\kopf}{%
  \noindent
  {\large\bfseries\color{accent} Numerische Methoden\, --\, Übungsklausur 01}\\[2pt]
  {\small\color{fgtext} Klausurnahe Aufgaben im Stil der Probeklausur \;·\; Open Book, kein Taschenrechner}\\[6pt]
  \noindent\textcolor{rulecol}{\rule{\linewidth}{0.4pt}}\\[4pt]
  {\small Name:\; \textcolor{rulecol}{\rule{6cm}{0.3pt}} \hfill Datum:\; \textcolor{rulecol}{\rule{4cm}{0.3pt}}}\\[2pt]
  \noindent\textcolor{rulecol}{\rule{\linewidth}{0.4pt}}
  \par\vspace{8pt}
}

% ----- Aufgabenkopf: #1 Nr/Titel, #2 Punkte, #3 Schwierigkeit -----
\newcommand{\aufg}[3]{%
  \needspace{3\baselineskip}
  \par\vspace{2pt}
  {\large\bfseries\color{accent2} Aufgabe #1}\hfill
  {\small\color{fgtext}#2 Punkte \;·\; #3}\\[2pt]
  \noindent\textcolor{rulecol}{\rule{\linewidth}{0.3pt}}
  \par\vspace{6pt}
}

\usepackage{needspace}

\setlength{\parindent}{0pt}
\setlength{\parskip}{4pt}

\begin{document}

\kopf

Bearbeitungshinweise: Begründen Sie Ihre Antworten kurz; eine reine Zahl ohne
Rechenweg gibt keine volle Punktzahl. Ableitungen dürfen analytisch gebildet
werden. Brüche dürfen stehen bleiben. Nutzen Sie das Karo-Raster für Ihre
Rechnung.

% =====================================================================
\clearpage
\aufg{1 \;--\; Gleitkommasystem \& Auslöschung}{12}{mittel}

\textbf{(a)} Sie entwerfen ein Gleitkommasystem (floating-point system) mit
$1$ Vorzeichenbit und $5$ weiteren Bits, die Sie frei auf Mantisse (mantissa)
und Exponent (exponent) verteilen können. Diskutieren Sie den Trade-off:
Welche Eigenschaft verbessert sich, wenn Sie ein Bit der \emph{Mantisse}
zuweisen, welche, wenn Sie es dem \emph{Exponenten} geben? Argumentieren Sie
über Auflösung (resolution) und Wertebereich (dynamic range). \hfill\textit{(4 P)}

\karo{6cm}

\clearpage
\textbf{(b)} Fixieren Sie die Aufteilung nun auf $3$ Mantissenbits und
$2$ Exponentenbits mit Bias $1$ (normalisierte Mantisse $1{.}b_1b_2b_3$ mit
implizitem führenden Bit). Bestimmen Sie das Maschinenepsilon (machine epsilon)
$\varepsilon_{\mathrm{mach}}$ dieses Systems. \hfill\textit{(3 P)}

\karo{5cm}

\clearpage
\textbf{(c)} Berechnen Sie in diesem System das Ergebnis des Ausdrucks
\[
  \frac{1}{(1 + 0{,}05) - 1}.
\]
Geben Sie den Nenner \emph{nach Rundung} sowie das Endergebnis an.
Vergleichen Sie mit dem mathematischen Wert und erklären Sie, warum hier die
Auslöschung (catastrophic cancellation) das Resultat zerstört. \hfill\textit{(5 P)}

\karo{8cm}

% =====================================================================
\clearpage
\aufg{2 \;--\; Festkomma, Rundung \& Fehlerarten}{9}{leicht}

\textbf{(a)} Ein Festkommaformat (fixed-point) reserviert ganzzahlige Bits und
$f$ Nachkommabits (Auflösung $2^{-f}$). Wie viele Nachkommabits $f$ benötigen
Sie mindestens, um $3{,}625$ \emph{exakt} darzustellen? Geben Sie die
Binärdarstellung an. Begründen Sie, warum $\tfrac13$ in \emph{keinem} solchen
Format exakt darstellbar ist. \hfill\textit{(3 P)}

\karo{5cm}

\clearpage
\textbf{(b)} Runden Sie $\tfrac{10}{3}$ auf $2$ Nachkommastellen (dezimal).
Geben Sie den \emph{absoluten} und den \emph{relativen} Rundungsfehler an
(als Bruch genügt). \hfill\textit{(3 P)}

\karo{5cm}

\clearpage
\textbf{(c)} Ordnen Sie jede der folgenden Fehlerquellen genau einer Kategorie
zu -- \emph{Modellfehler}, \emph{Abschneide-/Diskretisierungsfehler} oder
\emph{Rundungsfehler} -- und begründen Sie kurz:
\begin{enumerate}
  \item Eine Ableitung wird durch den Differenzenquotienten
        $\frac{f(x+h)-f(x)}{h}$ ersetzt.
  \item $\tfrac13$ wird im Rechner als $0{,}333$ gespeichert.
  \item In einer Fallsimulation wird der Luftwiderstand vernachlässigt.
\end{enumerate}
Diskretisieren Sie außerdem das Intervall $[0,2]$ mit $N=4$ Teilintervallen:
Geben Sie die Schrittweite $h=\tfrac{b-a}{N}$ und die Gitterpunkte an.
\hfill\textit{(3 P)}

\karo{6cm}

% =====================================================================
\clearpage
\aufg{3 \;--\; Gitter-Suche für ein lineares System}{8}{mittel}

Gesucht ist die Lösung des linearen Systems
\[
  \begin{pmatrix} 2 & 1 \\ 1 & 3 \end{pmatrix}
  \begin{pmatrix} w \\ b \end{pmatrix}
  =
  \begin{pmatrix} 3 \\ 5 \end{pmatrix}
\]
mit einer Gitter-Suche (grid search) über $w,b \in \{0,1,2\}$.

\textbf{(a)} Wählen Sie ein Fehlermaß (error) für die Gitter-Suche und
begründen Sie Ihre Wahl kurz. \hfill\textit{(2 P)}

\textbf{(b)} Bestimmen Sie anhand des gewählten Fehlermaßes die optimalen Werte
$w^\ast$ und $b^\ast$. \hfill\textit{(3 P)}

\textbf{(c)} Vergleichen Sie mit der exakten Lösung. Erklärt die Gitter-Suche
das System exakt? Wie könnten Sie die Gitter-Suche verbessern, und ist es mit
dieser Methode grundsätzlich möglich, die exakte Lösung zu treffen?
\hfill\textit{(3 P)}

\textit{Hinweis. Die exakte Lösung lautet $w^\star=\tfrac45$, $b^\star=\tfrac75$.}

\karo{9cm}

% =====================================================================
\clearpage
\aufg{4 \;--\; Fehleranalyse: Kondition, Stabilität, Konsistenz, Konvergenz}{10}{schwer}

Wir suchen das Minimum von $g(x) = (x-2)^2$ auf $[0,4]$ mit stückweise
konstanter, gitterbasierter Diskretisierung (piecewise constant discretization)
bei zwei Schrittweiten $h_1 = 1$ und $h_2 = \tfrac12$. Die diskretisierte
Funktion $\hat g$ wird am nächstgelegenen Gitterpunkt ausgewertet. Die Eingabe
sei durch Messfehler um $\tilde x = x \pm \tfrac12$ gestört.

\textbf{(a)} \emph{Kondition.} Quantifizieren Sie die Konditionierung des
Problems an mindestens zwei Stellen Ihrer Wahl und zeigen Sie den Unterschied
zwischen gut und schlecht konditioniertem Bereich. \hfill\textit{(3 P)}

\textbf{(b)} \emph{Stabilität.} Berechnen Sie die Stabilität von $\hat g$ an
einem Referenzpunkt Ihrer Wahl. Welche Schrittweite ist stabiler? \hfill\textit{(2 P)}

\textbf{(c)} \emph{Konsistenz.} Quantifizieren Sie die Konsistenz von $\hat g$
für $h_1$ und $h_2$. Welche Schrittweite ist konsistenter? \hfill\textit{(2 P)}

\textbf{(d)} \emph{Konvergenz.} Verbinden Sie Stabilität und Konsistenz zur
Konvergenz nahe dem Minimum. Führt eine kleinere Schrittweite \emph{zwingend}
zu einer besseren Approximation? Begründen Sie. \hfill\textit{(3 P)}

\karo{10cm}

% =====================================================================
\clearpage
\aufg{5 \;--\; Newton- \& Sekantenverfahren}{10}{schwer}

Gesucht ist eine Nullstelle (root) von
\[
  f(x) = x^3 - 2x + 2 .
\]

\textbf{(a)} Führen Sie ausgehend von $x_0 = 0$ zwei Iterationen des
Newton-Verfahrens (Newton's method) durch. $x_1$ und $x_2$ dürfen als Bruch
angegeben werden. \hfill\textit{(3 P)}

\textbf{(b)} Diskutieren Sie Ihre Beobachtung. Konvergiert das Verfahren? In
welchem weiteren Punkt würden Sie eine ähnliche Beobachtung machen? Nennen Sie
das Konvergenzkriterium (convergence criterion) des Newton-Verfahrens.
\hfill\textit{(3 P)}

\textbf{(c)} Nehmen Sie nun an, $f$ sei nur durch Funktionsauswertungen
zugänglich. Geben Sie eine numerische Approximation für $f'(x_0)$ an und nennen
Sie ihren Approximationsfehler (Ordnung in $h$). Leiten Sie daraus ein
ableitungsfreies Verfahren ab, mit dem der beobachtete Versagensmodus umgangen
werden kann. \hfill\textit{(4 P)}

\karo{9cm}

% =====================================================================
\clearpage
\aufg{6 \;--\; Globale Optimierung \& raue Verlustlandschaft}{11}{mittel}

\textbf{(a)} Betrachten Sie die Verlustlandschaft (loss landscape)
\[
  f(x) = (x-1)^2 + \tfrac12 \sin(4\pi x), \qquad x \in [0,2].
\]
Erklären Sie, warum der hochfrequente Sinus-Term für lokale Optimierer
(Gradientenabstieg, Newton) problematisch ist. \hfill\textit{(3 P)}

\textbf{(b)} Neben globalen Optimierungsmethoden gibt es einen weiteren Ansatz,
um die Oszillation zu umgehen. Nennen Sie diesen Ansatz und wählen Sie eine
geeignete Schrittweite $h$. Begründen Sie über die Periode des Sinus-Terms.
\hfill\textit{(3 P)}

\textbf{(c)} \emph{Random Restarts.} Ein zufälliger Neustart findet das globale
Minimum mit Wahrscheinlichkeit $p=\tfrac12$. Geben Sie die Erfolgs\-wahr\-schein\-lich\-keit
$P$ nach $K$ unabhängigen Neustarts an und bestimmen Sie das kleinste $K$ mit
$P \ge 0{,}99$. \hfill\textit{(3 P)}

\textbf{(d)} \emph{Simulated Annealing.} Erklären Sie kurz die Akzeptanzregel
$P_{\text{akz}} = e^{-\Delta E / T}$ und welche Rolle die sinkende Temperatur
$T$ spielt, um lokale Minima zu verlassen. \hfill\textit{(2 P)}

\karo{8cm}

% =====================================================================
\clearpage
\aufg{7 \;--\; Anfangswertproblem: Euler \& Runge-Kutta}{10}{mittel}

Gegeben sei das Anfangswertproblem
\[
  y'(t) = y(t), \qquad y(0) = 1, \qquad \text{exakte Lösung } y(t) = e^{t}.
\]
Verwenden Sie die Schrittweite $h = 1$ und führen Sie \emph{einen} Schritt durch.

\textbf{(a)} Berechnen Sie $y_1$ mit dem expliziten Euler-Verfahren
$y_{n+1} = y_n + h\,f(t_n,y_n)$ und geben Sie den Fehler zu $e \approx 2{,}718$
an. \hfill\textit{(3 P)}

\textbf{(b)} Berechnen Sie $y_1$ mit dem klassischen Runge-Kutta-Verfahren
RK4,
\[
  y_{n+1} = y_n + \tfrac{h}{6}\,(k_1 + 2k_2 + 2k_3 + k_4),
\]
inklusive aller $k_i$. Vergleichen Sie den Fehler mit (a). \hfill\textit{(4 P)}

\textbf{(c)} Geben Sie das Butcher-Tableau von RK4 an und nennen Sie die
Konvergenzordnung von Euler und RK4. \hfill\textit{(3 P)}

\karo{9cm}

% =====================================================================
\clearpage
\aufg{8 \;--\; Numerische Integration (Newton-Cotes)}{10}{mittel}

Betrachten Sie das Integral $\displaystyle \int_0^2 x^2 \,dx$ (exakter Wert
$\tfrac83$).

\textbf{(a)} Berechnen Sie die Näherung mit der \emph{Trapezregel} -- einmal mit
einem Intervall und einmal zusammengesetzt mit zwei Teilintervallen ($h=1$).
Geben Sie jeweils den Fehler an. \hfill\textit{(3 P)}

\textbf{(b)} Berechnen Sie die Näherung mit der \emph{Simpson-Regel}
($h=1$, zwei Teilintervalle). Erklären Sie, warum das Ergebnis exakt ist.
\hfill\textit{(3 P)}

\textbf{(c)} Berechnen Sie $\displaystyle \int_0^3 x^2 \,dx$ mit der
\emph{Simpson-$3/8$-Regel} (Stützstellen $0,1,2,3$). Notieren Sie die Gewichte.
\hfill\textit{(4 P)}

\karo{9cm}

% =====================================================================
\clearpage
\aufg{9 \;--\; Lagrange-Interpolation}{8}{mittel}

Gegeben seien die Stützpunkte $(0,1),\ (1,2),\ (2,5)$.

\textbf{(a)} Stellen Sie die Lagrange-Basispolynome $L_0, L_1, L_2$ auf.
\hfill\textit{(3 P)}

\textbf{(b)} Bestimmen Sie das Interpolationspolynom $p(x)$ und vereinfachen Sie
es zu einer Standardform. \hfill\textit{(3 P)}

\textbf{(c)} Werten Sie $p$ an der Stelle $x=\tfrac12$ aus. Nennen Sie ein
Problem hoher Polynomgrade (Runge-Phänomen) und eine Alternative.
\hfill\textit{(2 P)}

\karo{9cm}

% =====================================================================
\clearpage
\aufg{10 \;--\; Monte-Carlo-Integration}{8}{mittel}

Geschätzt werden soll $\displaystyle \int_0^1 x^2\,dx$ (exakt $\tfrac13$) per
Monte-Carlo, $\;I \approx |\Omega|\,\tfrac1N \sum_{i=1}^N f(X_i)$.

\textbf{(a)} Berechnen Sie die Schätzung für die $N=4$ Stützstellen
$X_i = 0{,}1;\ 0{,}3;\ 0{,}6;\ 0{,}9$ (mit $|\Omega| = 1$). Vergleichen Sie mit
dem exakten Wert. \hfill\textit{(3 P)}

\textbf{(b)} Der Standardfehler ist $\mathrm{SE} = |\Omega|\,\sigma_f/\sqrt{N}$.
Wie viele Stützstellen brauchen Sie etwa, um den Fehler zu \emph{halbieren}?
\hfill\textit{(2 P)}

\textbf{(c)} Erklären Sie den \emph{Fluch der Dimensionen} (curse of
dimensionality) und warum Monte-Carlo in hohen Dimensionen einem regelmäßigen
Gitter überlegen ist. \hfill\textit{(3 P)}

\karo{8cm}

% =====================================================================
\clearpage
\aufg{11 \;--\; Fourier-Analyse \& DFT/FFT}{9}{schwer}

\textbf{(a)} Ein Signal sei $f(t) = 3\cos(2\pi\cdot 1\cdot t) + \cos(2\pi\cdot 4\cdot t)$.
Welche Frequenzen und Amplituden erwarten Sie im Spektrum
$F(\omega)=\int f(t)e^{-i\omega t}\,dt$? Nutzen Sie
$\cos(\omega t) = \tfrac12(e^{i\omega t}+e^{-i\omega t})$. \hfill\textit{(3 P)}

\textbf{(b)} Berechnen Sie die DFT $X_k = \sum_{n=0}^{N-1} x_n\,e^{-i 2\pi k n/N}$
der Folge $x = (1,0,1,0)$, $N=4$, für $k=0,1,2,3$. Interpretieren Sie das
Spektrum. \hfill\textit{(4 P)}

\textbf{(c)} Geben Sie die Komplexität von DFT und FFT an und begründen Sie,
warum die FFT für große $N$ entscheidend ist. \hfill\textit{(2 P)}

\karo{8cm}

\end{document}
